\chapter*{ओपन लॉजिक प्रोजेक्ट के बारे में}
\addcontentsline{toc}{chapter}{ओपन लॉजिक प्रोजेक्ट के बारे में}


\textit{Open Logic Text}, औपचारिक अधितर्कशास्त्र और औपचारिक विधियों की
मुक्त-स्रोत सहयोगात्मक पाठ्यपुस्तक है, जो मध्यवर्ती स्तर (अर्थात् औपचारिक
तर्कशास्त्र के आरंभिक पाठ्यक्रम के बाद) से आरंभ होती है। यद्यपि इसका लक्ष्य
गैर-गणितीय पाठक (विशेषतः दर्शन और कंप्यूटर विज्ञान के विद्यार्थी) हैं, यह
कठोर है।

अभी सम्मिलित कुछ विषयों का विवेचन संभवतः पूरा नहीं है और अनेक खंडों में अब
भी पर्याप्त संशोधन चाहिए। भविष्य में अधिक विषयों को समेटने के लिए पाठ का
विस्तार करने की योजना है। हम इसमें शब्दावली, आगे पढ़ने की सूची, ऐतिहासिक
टिप्पणियाँ, चित्र, बेहतर व्याख्याएँ, दर्शन, कंप्यूटर विज्ञान और गणित के लिए
परिणामों की प्रासंगिकता समझाने वाले खंड तथा अधिक समस्याएँ और उदाहरण जैसी
विशेषताएँ जोड़ने की भी योजना रखते हैं। यदि आपको कोई त्रुटि मिले या कोई सुझाव
हो, तो \href{https://github.com/OpenLogicProject/OpenLogic/wiki/Contributing}{परियोजना
दल को बताएँ}।

परियोजना मुक्त स्रोत की भावना से चलती है। पाठ केवल निःशुल्क उपलब्ध ही नहीं
है; हम लाटेक स्रोत को क्रिएटिव कॉमन्स एट्रिब्यूशन लाइसेंस के अंतर्गत देते
हैं, जो उपयुक्त श्रेय देने की शर्त पर किसी को भी हमारा कार्य डाउनलोड करने,
प्रयोग करने, बदलने, पुनर्व्यवस्थित करने, रूपांतरित करने और पुनर्वितरित करने का
अधिकार देता है। अधिक जानकारी के लिए ओपन लॉजिक प्रोजेक्ट की वेबसाइट
\href{http://openlogicproject.org/}{openlogicproject.org} देखें।
